\documentclass[11pt]{article}
\usepackage[a4paper,margin=27mm]{geometry}
\usepackage{amsmath,amssymb,mathtools}
\usepackage{fontspec,unicode-math}
\setmainfont{Latin Modern Roman}
\setmathfont{Latin Modern Math}
\usepackage{xurl,hyperref}
\hypersetup{colorlinks=true,urlcolor=blue,linkcolor=blue}
\newcommand{\sref}[2]{\hyperlink{source:#1:#2}{[S#2]}}
\newcommand{\cataloguescope}[1]{\paragraph{Mathematical question.} #1}
\begin{document}
\section{Conjecture 2.15 — Mahler}\label{statement}

\subsection{Definitions and mathematical statement}

\paragraph{Shared notation.}$\mathbb N=\{1,2,\ldots\}$, and $\mathbb Z,\mathbb Q,\mathbb R,\mathbb C$ denote the integers, rationals, reals and complex numbers. Any other conventions are specified locally.


Let $(\varepsilon_n)_{n\ge0}$ be any sequence with $\varepsilon_n\in\{0,1\}$. The absolutely convergent series $\xi=\sum_{n\ge0}\varepsilon_n3^{-n}$ is real. A number is transcendental if it is not algebraic over $\mathbb Q$.
\paragraph{Statement recorded in the supplied source.}
For every such sequence, does $\xi\notin\mathbb Q$ imply that $\xi$ is transcendental? \sref{statement}{2}
\subsection{Short English statement}
Must every irrational number formed from only zero and one ternary digits be transcendental?
\subsection{Sources}
\begin{description}
\item[\hypertarget{source:statement:1}{S1}] ULAM.AI and the UnsolvedMath Contributors, UnsolvedMath: A Curated Collection of Open Mathematics Problems, catalogue entry. CC BY 4.0. \url{https://huggingface.co/datasets/ulamai/UnsolvedMath}.
\item[\hypertarget{source:statement:2}{S2}] Michel Waldschmidt, Open Diophantine Problems, arXiv:math/0312440v2 (24 January 2004), Conjecture 2.15. \url{https://arxiv.org/abs/math/0312440}.
\end{description}
\label{end:statement}
\end{document}
